\section{Related Work}

A mature market of commercial software-as-a-service platforms now supports remote usability evaluation. UserTesting combines a proprietary participant panel with video-first moderated (``Live Conversation'') and unmoderated studies overlaid with AI-assisted analysis~\cite{usertesting2026platform}, and has absorbed the quantitative benchmarking platform UserZoom following their 2023 merger~\cite{usertesting2026userzoom}. Maze is unmoderated by default, converting prototypes and live websites into task-based studies scored through automatically instrumented metrics such as completion rate, misclick rate, and task duration~\cite{maze2026platform}. Lookback specializes in live moderated sessions with screen, camera, and audio capture, stakeholder observation rooms, and timestamped free-form notes~\cite{lookback2026platform}; Loop11 supports both modes and, unusually, auto-computes SUS and NPS for unmoderated tests~\cite{loop112026platform}; Useberry targets prototype-centric unmoderated testing with click and flow analytics~\cite{useberry2026platform}. Across these tools, quantitative evidence derives almost entirely from automated client-side instrumentation in unmoderated studies, while moderated modes yield video recordings and free-form notes that must be coded after the fact. To the best of our knowledge, none \emph{combines} structured real-time evaluator logging in moderated sessions, per-task ease ratings integrated into the same pipeline, and an open, self-hostable architecture in one platform.

Academic work on remote evaluation dates to Hartson et al., who framed the network as an extension of the usability laboratory~\cite{hartson1996remote}. Andreasen et al.\ subsequently found remote synchronous testing ``virtually equivalent'' to conventional lab-based think-aloud, whereas asynchronous methods identify fewer problems at lower moderator cost~\cite{andreasen2007remote}, an asymmetry confirmed by Bruun et al.~\cite{bruun2009asynchronous} and persisting in recent comparisons~\cite{alhadreti2022comparison}. On the tooling side, RemUSINE and WebRemUSINE pioneered task-level instrumentation by comparing interaction logs against task models, but operate post hoc rather than steering live sessions~\cite{paterno2000remusine,paganelli2003webremusine}. TechSmith's Morae became the de facto observation suite in usability practice~\cite{techsmith2004morae,bastien2010usability}; its Observer component supported real-time marker and error logging by observers during moderated sessions and is the closest precedent to our evaluator instrument, but Morae is a discontinued desktop product requiring installed software on the participant's machine, without participant-side questioning or an open architecture. More recent browser-based platforms such as UTAssistant~\cite{desolda2017utassistant,federici2018utassistant} and the analytics-augmented test-management system of Generosi et al.~\cite{generosi2022test} target largely asynchronous workflows, and an action-research deployment of remote testing found the evaluation tool itself to be an obstacle to high-quality results~\cite{pedersen2021asynchronous}.

Avalux's measurement model follows the ISO 9241-11 decomposition of usability into effectiveness, efficiency, and satisfaction~\cite{iso2018usability}. Perceived usability is captured with the ten-item SUS~\cite{brooke1996sus}, whose reliability~\cite{bangor2008empirical,lewis2018sus}, benchmark mean of 68~\cite{sauro2011practical}, and adjective-based interpretation~\cite{bangor2009determining} are well established; per-task difficulty uses the Single Ease Question, shown to perform comparably to more elaborate post-task instruments~\cite{sauro2009comparison}. Completion rate, time-on-task, and action counts relative to an optimal path are the standard objective performance metrics~\cite{albert2022measuring,sauro2016quantifying}, and moderated testing with small samples is supported by classic problem-discovery models~\cite{virzi1992refining,nielsen1993mathematical}.

A final thread concerns where sessions take place and on whose device. Early comparisons found no significant remote-versus-local differences in critical errors or time-on-task~\cite{mcfadden2002remote,brush2004comparison}, while lab-versus-field studies of mobile applications show that removing participants from their natural device context can mask context-dependent problems~\cite{duh2006usability,kaikkonen2005usability}. The COVID-19 pandemic pushed moderated protocols onto general-purpose videoconferencing~\cite{simonliedtke2021remote,hill2021usability}, with studies of remote moderation identifying the moderator's reduced visibility into the participant's device and context as its principal cost~\cite{lobchuk2022remote}.

Avalux addresses the gap these threads leave open. It is an open, integrated platform for moderated testing in which the evaluator live-logs, per task, completion time, action counts, typed errors, hesitation events, and SEQ ratings, while the participant completes tasks on their own device via a browser join link, with task state synchronized between the two roles in real time over WebSockets. A single pipeline carries a session from protocol definition through per-task SEQ, post-session SUS, and a semi-structured interview to an automatically generated PDF report. Unlike the commercial services surveyed above, Avalux is open and self-hostable---a React single-page application over Supabase Postgres with row-level security and realtime channels---and unlike prior academic tooling it makes fine-grained evaluator-side logging and evaluator--participant realtime synchronization first-class features of the moderated session rather than post hoc analyses.
